\documentclass[12pt]{article}
\usepackage[utf8]{inputenc}
\usepackage[T1]{fontenc}
\usepackage{hyperref}
\title{Arabic, Judeo-Tripolitanian Language Resources for K-12 Curriculum Development}
\author{Compiled by Claude Sonnet 5.5 (claude-sonnet-5-5) for Andr\'as Kornai}
\date{}
\begin{document}
\maketitle

\section{Introduction}
Judeo-Tripolitanian Arabic (yud) was the Arabic variety of the Jews of Tripoli, Libya. Its speakers are now mostly elderly and bilingual in Hebrew, in Israel and Italy. This report rests on a few search results. No school materials were found.

\section{Primary Educational Resources}
No primers, textbooks or school programmes were found in the search results.

\section{Linguistic and Reference Materials}
Sumikazu Yoda, \emph{The Arabic Dialect of the Jews of Tripoli (Libya): grammar, text and glossary} (2005): \url{https://ub16.uni-tuebingen.de/Record/487455495}.
Luca D'Anna, \emph{Libyan Judeo-Arabic} (Semitica Viva 69, 2025): \url{https://www.harrassowitz-verlag.de/Libyan_Judeo-Arabic/title_9129.ahtml}, with a review at \url{https://silphiumgatherer.com/2026/01/26/book-libyan-judeo-arabic/}.
Overview: \url{https://en.wikipedia.org/wiki/Judeo-Tripolitanian_Arabic}.

\section{Religious and Community Resources}
Background on Libyan Jewish memory and identity: \url{https://mail.sephardichorizons.org/Volume8/Issue1&amp;2/Lagnado.html}.

\section{Digital and Audio Resources}
None were found in the search results.

\section{Curriculum Development Context}
A heritage or revitalization curriculum would have to be built from the scholarly grammar, texts and glossary above, together with diaspora communities. No existing curriculum was found.

\end{document}